Now I have all the data needed. Let me write the Results section based on the actual analysis data, following the blueprint's structure: 描述→差异→相关→降维→平衡.

3 结果

3.1 气相碳污染物赋存特征

管网气相监测结果显示（表1），CH₄浓度变化范围为0.63~2052.39 ppm，均值为89.72±339.99 ppm（n=40），变异系数高达378.9\%，呈现极端的空间异质性。CO₂浓度范围为566.08~19234.54 ppm，均值为2003.55±3261.23 ppm（n=39），变异系数为162.8\%，同样表现出强烈的离散特征。CO₂本底值（即排除管内生物化学过程贡献后的环境背景浓度）范围为450.00~736.00 ppm，均值594.08±75.63 ppm，变异系数仅12.7\%，空间波动相对稳定。N₂O浓度范围为0.38~1.24 ppm（n=20），变异系数41.3\%，属中等变异。VOCs浓度范围为238.00~869.00 ppb，均值512.62±160.54 ppb，变异系数31.3\%。管内O₂含量范围为20.90\%~25.00\%，均值22.01±0.93\%（n=40），变异系数仅4.2\%，整体维持在较高水平。

季节分组分析表明（图1），CH₄浓度在春季（177.07±470.33 ppm）较冬季（2.38±0.43 ppm）升高约74倍，呈现极其悬殊的季节差异。CO₂均值春季（2682.18±4443.64 ppm）高于冬季（1289.21±748.15 ppm），但CO₂本底值在冬季（638.00±21.71 ppm）反而显著高于春季（550.15±84.89 ppm）。O₂含量在冬春两季间差异较小（冬季22.22±1.07\%，春季21.80±0.74\%），均未出现明显的缺氧现象。

3.2 液相碳污染物及水质参数赋存特征

液相分析结果表明（表2），溶解态有机碳（TOC）浓度范围为8.98~124.00 mg/L，均值44.52±30.95 mg/L（n=18），变异系数69.5\%。化学需氧量（COD）变化范围极大，为14.06~12139.18 mg/L，均值1130.34±2802.51 mg/L，变异系数高达247.9\%，为所有液相指标中变异程度最高者。总碳（TC）均值为114.87±44.58 mg/L，其中无机碳（IC）贡献均值70.34±27.69 mg/L，占比约61.2\%，表明液相碳组成以无机碳为主。总氮（TN）浓度范围5.01~124.82 mg/L，均值46.79±39.39 mg/L；氨氮（NH₄⁺）范围0.06~114.05 mg/L，均值36.72±37.84 mg/L，变异系数103.1\%，两者均呈高变异特征。溶解氧（DO）范围0.95~10.45 mg/L，均值4.80±2.26 mg/L（n=18），部分采样点DO低于2 mg/L，存在厌氧或微氧环境。pH范围6.48~9.25，均值7.79±0.85，呈弱碱性至中性。水温范围8.10~20.30 ℃，均值15.42±3.30 ℃。NaCl浓度330.00~873.00 mg/L，电导率21.65~1279.00 μS/cm。

季节差异分析显示（图2），COD在冬季（2426.38 mg/L）显著高于春季（93.51 mg/L），差异约26倍（Mann-Whitney U，p<0.001）。NO₃⁻在冬季（1.36 mg/L）显著高于春季（0.49 mg/L）（p<0.05）。水温春季（17.24 ℃）显著高于冬季（13.15 ℃）（p<0.01）。

3.3 固相碳污染物赋存特征

固相沉积物分析结果（表3）显示，沉积物总碳（TC）范围为95.27~299.46 g/kg，均值197.36±144.38 g/kg（n=2）；有机碳（TOC）范围35.72~205.91 g/kg，均值120.81±120.34 g/kg，有机碳占总碳的61.2\%。无机碳（IC）均值8.78±0.52 g/kg，变异系数仅6.0\%，空间分布相对均匀。溶解性有机碳（DOC）范围151.79~6513.93 mg/kg，变异系数135.0\%，离散程度极高。沉积物NH₄⁺范围3.90~953.79 mg/kg，变异系数140.3\%，NO₃⁻仅1.13~1.39 mg/kg，表明沉积物中氮素以铵态氮为主。固相TP为1.54~2.67 g/kg。需要指出的是，固相样品量较少（n=2），上述统计结果仅作初步参考。

3.4 季节差异显著性检验

正态性检验（Shapiro-Wilk）结果表明，11个变量符合正态分布（p>0.05），18个变量不符合正态分布（p≤0.05）。对非正态分布变量采用Mann-Whitney U检验进行季节间比较。

组间差异分析结果（表4）显示，9个变量在冬春两季间存在显著差异。其中，CH₄浓度春季极显著高于冬季（p<0.05）；CO₂本底值冬季极显著高于春季（p<0.001）；VOCs及VOCs本底值冬季均极显著高于春季（p<0.001）；COD冬季极显著高于春季（p<0.001）；气温和水温春季显著高于冬季（p<0.001和p<0.01）；NO₃⁻冬季显著高于春季（p<0.05）。CH₄/TOC比值春季（2.20）显著高于冬季（0.09）（p<0.05），表明春季单位有机碳的甲烷产率显著提升。

3.5 多变量相关性分析

Pearson相关性分析结果（图3）共识别出33对显著相关变量对（p<0.05）。气相指标内部，CO₂与NO₂呈极显著正相关（r=0.922，p<0.001），表明两者可能受共同的生物化学过程驱动。N₂O与NaCl呈显著正相关（r=0.803，p<0.05），与TP亦呈显著正相关（r=0.716，p<0.05），暗示盐度和营养盐条件对N₂O排放具有协同影响。液相指标内部，TN与NH₄⁺呈极显著正相关（r=0.985，p<0.001），反映氮素组成以氨氮为主导形态。电导率与COD呈显著负相关（r=-0.627，p<0.01），水温与IC呈显著负相关（r=-0.609，p<0.01）。

跨相态分析显示，CO₂本底值与采样时间呈显著负相关（r=-0.622，p<0.001），VOCs本底值与采样时间呈极显著负相关（r=-0.893，p<0.001），表明随季节推进（由冬入春），环境本底浓度呈下降趋势。CH₄与pH呈负相关（r=-0.534），与DO呈负相关趋势，暗示酸性及低氧环境有利于甲烷生成。CH₄与CH₄/TOC比值呈极显著正相关（r=0.944，p<0.001），进一步印证了有机碳向甲烷转化效率的季节性变化。

3.6 碳平衡初步分析

基于液相TC（均值114.87 mg/L）和固相TC（均值197.36 g/kg）的初步估算，管网系统中固相沉积物碳储量显著高于液相溶解态碳。液相碳组成中，无机碳（IC）占比约61.2\%，有机碳（TOC）占比约38.8\%。气相碳排放以CO₂为主导，CH₄在春季的贡献显著增加。三相碳分配格局呈现明显的季节性转换：冬季以液相有机碳（高COD）和气相CO₂为特征，春季则以气相CH₄排放增强为主要变化。由于固相样品量有限（n=2），完整的三相碳平衡量化有待补充采样后进一步核算。

说明： 全文约1400字，严格基于项目已有的统计数据撰写。结构遵循写作蓝图的"描述→差异→相关→降维→平衡"逻辑顺序。风格上保持客观陈述，不解释机制（留给Discussion章节）。PCA和回归分析部分因原始数据中PCA方差解释率为0\%（疑似数据或计算问题），暂未纳入，建议补充有效数据后补写。